\subsection{Long-Range Dependence and Self-Similarity}

\paragraph{Leland, Taqqu, Willinger, and Wilson (1994).}
\emph{On the self-similar nature of Ethernet traffic.}
This paper is one of the central empirical starting points for modern LRD
network-traffic modelling. It showed that packet traffic could retain burstiness
across aggregation scales, motivating models beyond Poisson and short-memory
Markov assumptions. For S5.jl it supplies the broader methodological lesson:
synthetic test data should reproduce large-scale structure, not just local
statistics.

\paragraph{Pipiras and Taqqu (2017).}
\emph{Long-Range Dependence and Self-Similarity.}
This monograph provides the mathematical background tying together covariance
decay, spectral poles, self-similarity, and the Hurst parameter. It is the main
reference for the parameter relationships used throughout the package, including
the links between spectral exponent, ACF decay exponent, and Hurst parameter.

\paragraph{Beran (1994).}
\emph{Statistics for Long-Memory Processes.}
Beran's text is a core statistical reference for inference and modelling in
long-memory time series. Although S5.jl is not an estimator package, the book
helps frame why estimator benchmarking needs synthetic sequences with known
parameters and known violations of short-memory assumptions.

\paragraph{Taqqu, Teverovsky, and Willinger (1995).}
\emph{Estimators for long-range dependence: an empirical study.}
This empirical comparison of LRD estimators motivates the separation between
S5.jl and a future estimator package. The paper demonstrates that estimator
performance depends heavily on model class, sample size, and diagnostics, which
is why S5.jl focuses first on controllable generators.

\subsection{Numerical LRD Synthesis}

\paragraph{Paxson (1997).}
\emph{Fast, approximate synthesis of fractional Gaussian noise for generating
self-similar network traffic.}
This is the direct ancestor of S5.jl's PB1 method. It gives a practical spectral
construction of fGn, trading exact covariance matching for speed. In S5.jl the
numerical fGn sequence is mapped to symbols by rank binning, preserving a simple
path from target Hurst parameter to categorical output.

\paragraph{Dieker (2004).}
\emph{Simulation of fractional Brownian motion.}
Dieker's thesis surveys exact and approximate fBm/fGn simulation methods,
including circulant embedding. It provides the natural next reference for
improving PB1 if exact covariance control becomes more important than speed.

\paragraph{Roughan, Veitch, and Abry (2001).}
\emph{Real-time estimation of the parameters of long-range dependence.}
This work is important for the wavelet perspective behind the planned PB3
generator. Wavelet methods are attractive because they align naturally with
multi-scale structure and may provide a bridge from LRD synthesis to local
state-machine control.

\subsection{Point Processes and Heavy-Tailed Models}

\paragraph{Garrett and Willinger (1994).}
\emph{Analysis, modeling, and generation of self-similar VBR video traffic.}
This paper is a classic example of generating self-similar traffic from
heavy-tailed source behaviour. It motivates MB2's heavy-tailed regime-sojourn
construction, where long periods in a regime create large-scale dependence while
the within-regime Markov chain controls local symbol structure.

\paragraph{Lowen and Teich (1995).}
\emph{Estimation and simulation of fractal stochastic point processes.}
Lowen and Teich provide the point-process foundation for fractal renewal and
shot-noise constructions. Their work informs MB3's view of each symbol stream as
an independent heavy-tailed arrival process.

\paragraph{Ryu and Lowen (1998).}
\emph{Point process models for self-similar network traffic, with applications.}
This paper develops point-process models that connect heavy-tailed mechanisms to
self-similar traffic. It is a key bridge between network-traffic synthesis and
the symbolic FSS construction used in S5.jl.

\paragraph{Roughan, Yates, and Veitch (1999).}
\emph{The mystery of the missing scales: pitfalls in the use of fractal renewal
processes to simulate LRD processes.}
This paper is a warning label for naive heavy-tailed simulation. It shows that
FRP constructions can fail to express LRD over the intended scale range. S5.jl's
validation scripts and TODO list retain this as an explicit design concern for
FSS and On/Off methods.

\subsection{Symbolic and Categorical Models}

\paragraph{Gal, Chen, and Ghahramani (2015).}
\emph{Latent Gaussian processes for distribution estimation of multivariate
categorical data.}
This work motivates PB2. The central idea is to represent categorical outcomes
through latent Gaussian variables and an argmax transformation. S5.jl adapts that
idea to long-memory latent fGn streams and adds finite-sample offset calibration
for marginal control.

\paragraph{Kumar, Raghu, Sarlos, and Tomkins (2017).}
\emph{Linear additive Markov processes.}
This paper motivates MB1. LAMP models make transition probabilities additive in
past observations, allowing long-range memory to be encoded through a decay
kernel over history. S5.jl implements a finite-depth version with a marginal
innovation term.

\paragraph{Singh, Greenberg, and Klakow (2016).}
\emph{The custom decay language model for long range dependencies.}
This language-modelling paper is closely related to LAMP in spirit: it uses a
custom decay over historical context to model long-range dependencies in text.
It supports the idea that symbolic sequences need memory mechanisms distinct
from ordinary short-context Markov models.

\paragraph{Carpio and Daley (2007).}
\emph{Long-range dependence of Markov chains in discrete time on countable state
space.}
This paper is important because it frames LRD for Markov chains through counting
and recurrence-time behaviour. It helps justify the use of count-process and
return-time definitions for non-numerical sequences where ordinary numerical ACF
methods are not directly applicable.

\subsection{Machine Learning Motivation}

\paragraph{Belletti, Beutel, Jain, and Chi (2018).}
\emph{Factorized recurrent neural architectures for longer range dependence.}
This paper motivates the project from the machine-learning side by showing that
common recurrent architectures struggle with long-range dependence. It helps
explain why synthetic NN-LRD data is useful for stress-testing learning
algorithms.

\paragraph{Belletti, Chen, and Chi (2019).}
\emph{Quantifying long range dependence in language and user behavior to improve
RNNs.}
This paper is one of the direct motivations for studying LRD in language and
human behavioural sequences. It connects empirical LRD measurement to model
design, reinforcing S5.jl's role as a controllable data source for later
estimator and learning experiments.
